\documentclass[12pt]{article}
\usepackage[T2A]{fontenc}
\usepackage[utf8]{inputenc}
\usepackage[ukrainian]{babel}
\usepackage[hidelinks]{hyperref}
\usepackage{amsmath,amssymb,amsthm,thmtools}
\usepackage{mathtools}
\usepackage{tikz,tikz-cd}
\usetikzlibrary{positioning,decorations.markings}
\usepackage{subcaption}

\usepackage{booktabs}
\usepackage{subcaption}
\usepackage{mathtools}
\mathtoolsset{centercolon,mathic}
\usepackage{tikz-cd}
\usetikzlibrary{positioning,decorations.markings}

\pgfdeclarearrow{% An arrowhead that matches the one from tx fonts
  name = txto,
  setup code = {
    % The different end values:
    \pgfarrowssettipend{1.9285\pgflinewidth}
    \pgfarrowssetbackend{-2.125\pgflinewidth}
    \pgfarrowssetlineend{0pt}
    \pgfarrowssetvisualbackend{0pt}
    % The hull
    \pgfarrowshullpoint{1.9285\pgflinewidth}{0pt}
    \pgfarrowsupperhullpoint{-1.5714\pgflinewidth}{3.6607\pgflinewidth}
    \pgfarrowsupperhullpoint{-2.125\pgflinewidth}{3.0714\pgflinewidth}
    % No saves
  },
  drawing code = {
    \pgfsetdash{}{0pt}%
    \pgfpathmoveto{\pgfpoint{1.9285\pgflinewidth}{0pt}}%
    \pgfpathlineto{\pgfpoint{-1.5714\pgflinewidth}{3.6607\pgflinewidth}}%
    \pgfpathlineto{\pgfpoint{-2.125\pgflinewidth}{3.0714\pgflinewidth}}%
    \pgfpathlineto{\pgfpoint{0pt}{0.5\pgflinewidth}}%
    \pgfpathlineto{\pgfpoint{0pt}{-0.5\pgflinewidth}}%
    \pgfpathlineto{\pgfpoint{-2.125\pgflinewidth}{-3.0714\pgflinewidth}}%
    \pgfpathlineto{\pgfpoint{-1.5714\pgflinewidth}{-3.6607\pgflinewidth}}%
    \pgfpathclose%
    \pgfusepathqfill
  }
}

\tikzset{>=txto}
% This is a bit of a hack: we want the TikZ line
% widths to depend on the font size:
% use 0.56pt for 10pt text, 0.504pt for 9pt text
% (is there a cleaner way to do this?)
\newlength\myfontsize
\newlength\mythinwidth
\newlength\mythickwidth
\newlength\mydoubledist
\makeatletter
\setlength{\myfontsize}{\f@size pt}
\makeatother
\setlength{\mythinwidth}{0.056\myfontsize}
\setlength{\mythickwidth}{2\mythinwidth}
\setlength{\mydoubledist}{0.147\myfontsize}
\tikzset{every picture/.style={line width={\mythinwidth}}}
\tikzset{thin/.style={line width={\mythinwidth}}}
\tikzset{thick/.style={line width={\mythickwidth}}}
\tikzset{arrow/.style={->,node font=\small}}
\tikzcdset{arrow style=tikz}

\usepackage[hidelinks]{hyperref}

% MACROS %%%%%%%%%%%%%%%%%%%%%%%%%%%%%%%%%%%%%%%%%%%%%%%%%%%%%%%%
\newcommand*{\Z}{\mathbb{Z}}
\newcommand*{\N}{\mathbb{N}}
\newcommand*{\calC}{\mathcal{C}}
\newcommand*{\calD}{\mathcal{D}}
\newcommand*{\calS}{\mathcal{S}}
\newcommand*{\Ntwo}{\mathbb{N}_{-2}}
\newcommand*{\BZ}{\mathbf{B}\mathbb{Z}}
\newcommand*{\inftyone}{\ensuremath{(\infty,1)}}
\newcommand*{\inftyGpd}{\ensuremath{\infty\mathbf{Gpd}}}
\newcommand*{\SmoothInftyGpd}{\ensuremath{\mathbf{Smooth}\infty\mathbf{Gpd}}}
\newcommand*{\Type}{\ensuremath{\mathbf{Type}}}
\newcommand*{\sType}{\ensuremath{\mathbf{sType}}}
\newcommand*{\ssType}{\ensuremath{\mathbf{ssType}}}
\newcommand*{\Prop}{\ensuremath{\mathbf{Prop}}}
\newcommand*{\Set}{\ensuremath{\mathbf{Set}}}
\newcommand*{\sSet}{\ensuremath{\mathbf{sSet}}}
\newcommand*{\Gpd}{\ensuremath{\mathbf{Gpd}}}
\newcommand*{\cgTop}{\ensuremath{\mathbf{Top}_{\mathrm{cg}}}}
\newcommand*{\blank}{\mathord{\hspace{1pt}\text{--}\hspace{1pt}}}
\newcommand*\Fin[1]{\overline{#1}}
\mathchardef\texthyphen="2D
\newcommand*{\trType}[1]{\ensuremath{#1\mathbf{\texthyphen Type}}}
\newcommand*{\inftyoneCat}{\ensuremath{(\infty,1)\mathbf{\texthyphen Cat}}}
\newcommand*{\op}{^{\mathrm{op}}}
\newcommand*{\seq}{\overset{\mathrm{s}}{=}}
\newcommand*{\HoTTua}{HoTT$_{\mathrm{ua}}$}
\newcommand*{\HoTTel}{HoTT$_{\mathrm{el}}$}
\newcommand*{\HoTTelp}{HoTT$_{\mathrm{el}{+}{+}}$}
\newcommand*{\HoTTuf}{HoTT$_{\mathrm{UF}}$}
\DeclareMathOperator{\pt}{pt}
\DeclareMathOperator{\id}{id}
\DeclareMathOperator{\fst}{fst}
\DeclareMathOperator{\snd}{snd}
\DeclareMathOperator{\Id}{Id}
\DeclareMathOperator{\fib}{fib}
\DeclareMathOperator{\iso}{iso}
%\DeclareMathOperator{\Fin}{Fin}
\DeclareMathOperator{\isContr}{isContractible}
\DeclareMathOperator{\isProp}{isProp}
\DeclareMathOperator{\isSet}{isSet}
\DeclareMathOperator{\isGpd}{isGroupoid}
\DeclareMathOperator{\isEquiv}{isEquiv}
\DeclareMathOperator{\isIso}{isIso}
\DeclareMathOperator{\hasDimen}{hasDimension}
\DeclareMathOperator{\univalence}{univalence}
\DeclareMathOperator{\base}{base}
\DeclareMathOperator{\cloop}{loop}
\DeclareMathOperator{\inl}{left}
\DeclareMathOperator{\inr}{right}
\DeclareMathOperator{\cglue}{glue}
\DeclarePairedDelimiter\trunc{\lVert}{\rVert} % truncation
\DeclarePairedDelimiter\tr{\lvert}{\rvert} % introduction for truncation
\DeclarePairedDelimiter\abs{\lvert}{\rvert} % geometric realization
\DeclarePairedDelimiter\sem{[\![}{]\!]} % semantics

%%% Local Variables:
%%% mode: latex
%%% TeX-master: higher-structure.tex
%%% End:


\begin{document}

\title{Вищі структури в гомотопічній теорії типів}
\author{Ульрік Бухгольц \\ 
  \small Факультет математики, Технічний університет Дармштадта, \\
  \small Schlossgartenstra\ss e 7, D-64289 Дармштадт, \\
  \small \texttt{buchholtz@mathematik.tu-darmstadt.de}
}
\date{}

\maketitle

\begin{abstract}
Передбачуваною моделлю гомотопічних теорій типів, що використовуються в Унівалентних основах, є $\infty$-категорія гомотопічних типів, також відомих як $\infty$-групоїди. Проблема \emph{вищих структур} полягає в побудові гомотопічних типів, необхідних для математики, особливо тих, які не є множинами. Поточного репертуару конструкцій, включаючи звичайні конструктори типів та вищі індуктивні типи, достатньо для багатьох, але не для всіх з них. Ми обговорюємо проблемні випадки, зазвичай ті, що містять нескінченну ієрархію даних когерентності, такі як напівсимпліціальні типи, а також проблему розвитку метатеорії гомотопічних теорій типів в Унівалентних основах. Ми також обговорюємо деякі запропоновані рішення.
\end{abstract}

\section{Вступ}
\label{sec:introduction}

Гомотопічна теорія типів є водночас фундаментальним починанням, метою якого є забезпечення нових основ для математики, і галуззю математики та логіки, метою якої є забезпечення інструментів для математичного аналізу гомотопічних (вищих розмірних) структур. Назвімо перше Унівалентними основами (UF), а друге --- Гомотопічною теорією типів (HoTT), розуміючи, що HoTT охоплює багато різних конкретних теорій типів.

У цьому розділі ми використовуємо питання \emph{вищих структур} як призму, крізь яку вивчаємо обидві ці цілі та їхні зв'язки з іншими підходами до основ математики.

Щоб мотивувати проблему вищих структур, нам потрібно нагадати, що передбачуваний універсум UF та основна модель HoTT є царством $\infty$-групоїдів, гомотопічного типу алгебраїчних структур, які мають елементи, ототожнення, ототожнення між ототожненнями тощо до нескінченності (ad infinitum), і ці ототожнення поводяться розумно в тому сенсі, що ми можемо обертати їх, компонувати та здійснювати над ними операцію віскерінгу (whisker), але очікувані закони виконуються лише з точністю до вищих ототожнень. \emph{Гомотопічна гіпотеза} Гротендіка говорить нам, що $\infty$-групоїди є тим самим, що й \emph{гомотопічні типи}, тому ми використовуватимемо ці терміни як взаємозамінні, з невеликою перевагою до останнього, оскільки тоді гомотопічні теорії типів є водночас теоріями гомотопічних типів, а також гомотопічними теоріями типів.

Поширеною помилкою є те, що вищі гомотопічні типи зустрічаються лише в теорії гомотопій або мають відношення лише до неї. Це зовсім не так, оскільки навіть тип множин, як він використовується в більшості математичної практики, є $1$-типом. І вищі структури зараз займають помітне місце в багатьох галузях: від геометрії, алгебри та теорії чисел до математики квантової теорії поля у фізиці та паралельних обчислень в інформатиці. Вступ до гомотопічних типів та гомотопічної гіпотези подано в розділі~\ref{sec:inftygroupoids}.

Ключовою відмінністю між UF та більш ранніми підходами до основ, натхненними теорією категорій, є те, що UF приймає $\infty$-групоїди, а не різні поняття вищих категорій, за базові об'єкти математики, з яких решта отримуються шляхом додавання подальшої структури. Це розуміння належало Воєводському\footnote{В~\cite{Voevodsky2014} він писав: «Найбільшою перешкодою для мене була ідея, що категорії — це «множини в наступному вимірі». Я чітко пам’ятаю відчуття прориву, яке я відчув, коли зрозумів, що ця ідея хибна. Категорії не є «множинами в наступному вимірі». Вони є «частково впорядкованими множинами в наступному вимірі», а «множинами в наступному вимірі» є групоїди.»}, який зауважив, що багато природних конструкцій не є функторіальними в сенсі теорії категорій. (Подумайте, наприклад, про центр групи.) Проте кожна конструкція — якщо вона має математичний зміст — повинна зберігати відповідне поняття еквівалентності. (Ізоморфні групи дійсно мають ізоморфні центри тощо).

Оскільки UF прагне бути основою для всієї математики, необхідно, щоб її мова у формі HoTT забезпечувала засоби конструювання для всіх гомотопічних типів, які використовуються в математиці. Для побудови \emph{множин} це не така велика проблема, оскільки більшість множин, що зустрічаються в математиці, можуть бути побудовані за допомогою конструкторів типів теорії типів Мартіна-Льофа. (Але навіть тут є тонкощі, якщо ми хочемо залишатися в конструктивній і предикативній сфері.)

Основні проблеми виникають, коли йдеться про \emph{вищі розмірні} гомотопічні типи. Ми обговорюємо деякі позитивні результати (структури, які вже були побудовані), а також деякі відкриті проблеми (структури, які ще не були побудовані) в розділі~\ref{sec:problems}.

Зазначимо, що хоча ми очікуємо деяких фактичних \emph{негативних} результатів (тобто доведень неможливості) для деяких із відкритих проблем стосовно деяких конкретних гомотопічних теорій типів, вони ще не з'явилися. Але передбачаючи, що будуть потрібні подальші засоби побудови, ми обговорюємо потенційні рішення в розділі~\ref{sec:solutions}.

Для решти цього вступу ми розглянемо аналогію. Теорію типів Мартіна-Льофа можна розглядати як формальну систему для створення конструкцій. Фактично, варіант з імпредикативним універсумом був названий \emph{Численням конструкцій} (Calculus of Constructions, CoC), а подальше розширення, \emph{Числення індуктивних конструкцій} (Calculus of Inductive Constructions, CIC), є основою для доказового асистента Coq. І ми будемо стурбовані питанням меж методів конструювання, доступних у конструктивних теоріях типів. Очевидна аналогія напрошується сама собою, а саме з евклідовою геометрією та межами методів геометричних побудов за допомогою циркуля та лінійки. Ми (ймовірно) виявимо, що, як і в геометричному випадку, певні об'єкти не можна побудувати з найпростіших конструкцій, і вони вимагають для своєї побудови додаткових інструментів, таких як \emph{невсіс}. Однак ми будемо дотримуватися припису ощадливості Паппа і вимагатимемо, щоб усе, що можна побудувати меншими засобами, було побудовано саме так. Як наслідок, оскільки доведення є окремим випадком конструкції, ми вимагаємо, щоб якщо щось може бути доведено в слабшій системі, то воно має бути доведено саме так.

Очевидно, ми можемо включити до списку подальших засобів конструювання такі відомі принципи, як закон виключеного третього (LEM), принцип Маркова (MP), аксіому вибору (AC), різні види трансфінітної індукції (TI), а також принципи імпредикативності. Деякі з них, а також їх слабші версії, називаються \emph{конструктивними табу}, оскільки їх прийняття суперечить певним філософським поглядам, натхненним конструктивізмом або інтуїціонізмом, а також тому, що їх взагалі неможливо механічно виконати або це можливо лише зі значно збільшеною обчислювальною складністю.

Ще один аспект конструктивних табу полягає в тому, що вони зменшують кількість \emph{моделей}, в яких ми можемо інтерпретувати конструкції. Добре відомо, що конструктивні системи допускають багато корисних моделей, і справді, це одна з причин, чому класичні математики можуть бути зацікавлені в таких системах. Негомотопічні конструктивні системи часто можна моделювати в \emph{топосах}, точніше, $1$-топосах, які можна розглядати як узагальнення моделей Кріпке, як узагальнені простори, або дійсно як узагальнені світи множин.

Існує підозра, що HoTT можна моделювати у вищих топосах, точніше, \emph{\inftyone-топосах}. Вони різко розширюють корисність HoTT, як, наприклад, пояснюється в розділі Шрайбера. Більш ранні розширення теорії типів Мартіна-Льофа часто нав'язували аксіоми, такі як \emph{унікальність доведень тотожності} (UIP), які виключають вищі розмірні моделі. Вони суперечать аксіомі унівалентності і можуть бути названі \emph{гомотопічними табу}. Більш витончені аксіоми можуть виконуватися в \inftyone-топосах, що відповідають $1$-топосам ($1$-локальні \inftyone-топоси~\cite[Розділ~6.4]{LurieHTT}), але не в більш загальних \inftyone-топосах. Вони називаються \emph{конструктивно-гомотопічними табу}.

\section{Нескінченні групоїди та гомотопічна гіпотеза}
\label{sec:inftygroupoids}

Типи в UF передбачаються як \emph{гомотопічні типи}, тому зупинімося трохи на тому, чим вони є, як з інтуїтивної точки зору, так і з точки зору математики, розвиненої в теоретико-множинних основах.

Інтуїцію завжди важко передати, а у випадку з поняттям гомотопічного типу — тим більше. Інтуїція, зрештою, найкраще розвивається через практику та ознайомлення. Один із способів побудувати інтуїцію для гомотопічних типів полягає в роботі в гомотопічній теорії типів, або на папері, або за допомогою доказового асистента (proof assistant). Багато молодих дослідників у HoTT/UF робили це до вивчення теорії гомотопій з класичної точки зору.

У першому наближенні ми можемо сказати, що типи $A$ є сукупностями об'єктів разом із, для кожної пари об'єктів $a,b:A$, типом \emph{ототожнень} $p : a=_A b$, разом зі змістовними операціями над цими ототожненнями, такими як можливість їх компонувати та обертати. І також мають бути операції вищих порядків, які продукують ототожнення між ототожненнями, такі як ототожнення $\alpha(p) : (p^{-1})^{-1} = p$ для будь-якого $p : a = b$. Цей опис має на меті відобразити типи в їхньому втіленні як \emph{$\infty$-групоїди}, і з цієї точки зору два типи $A,B$ можна ототожнити, якщо існує (слабкий) функтор $F:A \to B$, який є еквівалентністю $\infty$-групоїдів.

Інша інтуїція походить від опису типів як (хороших) топологічних просторів з точністю до гомотопічної еквівалентності. Об'єкти — це точки простору, а ототожнення — це шляхи між точками. 

\emph{Гомотопічна гіпотеза} — це ідея про те, що ці окремі інтуїції відображають одне й те саме глибинне поняття. Вона виросла з гомотопічної гіпотези Гротендіка щодо конкретного визначення $\infty$-групоїдів~\cite{Grothendieck1983}. Сучасна термінологія належить Баєсу (Baez)~\cite{Baez2007}.

Щоб пояснити тонкість ситуації, звернімося до найпоширенішої реалізації ідеї $\infty$-групоїдів у контексті теоретико-множинної математики. Тут вони представлені \emph{симпліціальними множинами}, що задовольняють певну умову заповнення. Ці симпліціальні множини називаються \emph{комплексами Кана} на честь~\cite{KanIII}. Симпліціальна множина — це функтор $X : \Delta\op \to \Set$, де $\Delta$ — категорія непорожніх скінченних ординалів та функцій, що зберігають порядок. Це конкретно означає, що симпліціальна множина складається з множини $n$-симплексів $X_n$ для кожного $n=0,1,\dots$ разом із відображеннями граней та виродження, що задовольняють закони, які називаються \emph{симпліціальними тотожностями}. Ми мислимо $0$-симплекси як точки, $1$-симплекси як лінії між точками, $2$-симплекси як трикутники тощо.

Умова заповнення Кана стверджує, що якщо нам задано $n$ сумісних $(n-1)$-симплексів в $X$ у тому сенсі, що вони могли б бути $n$ з $n+1$ граней $n$-симплекса, то існує деякий такий $n$-симплекс. Ця умова проілюстрована на Рис.~\ref{fig:KanComplex} у деяких випадках малої розмірності. У кожному випадку ми можемо розглядати задані дані як відображення з \emph{рога} (horn), підсимпліціальної множини $\Lambda^n_k\subseteq\Delta^n$ стандартного $n$-симплекса $\Delta^n$, що складається з об'єднання всіх граней, протилежних $k$-й вершині, в $X$. Підняття (lift) — це деяке продовження цього до відображення з $\Delta^n$ в $X$, або, що еквівалентно, $n$-симплекс в $X$ із необхідними гранями.

\begin{figure}
  \centering
  \begin{subfigure}[b]{0.21\linewidth}
    \centering
    \begin{tikzpicture}[dot/.style={draw,shape=circle,
        minimum size=1pt,inner sep=0pt,outer sep=0pt,fill=black}]
      \coordinate[dot,label=210:$0$] (p0) at (0,0);
      \coordinate[dot,label=$1$] (p1) at (60:1);
      \coordinate[dot,label=-30:$2$] (p2) at (1,0);
      \begin{scope}[decoration={markings,
          mark=at position 0.55 with {\arrow{>}}}]
        \draw[postaction={decorate}] (p0) -- (p1);
        \draw[postaction={decorate}] (p0) -- (p2);
        \draw[dashed,postaction={decorate}] (p1) -- (p2);
      \end{scope}
    \end{tikzpicture}
    \caption{$\Lambda^2_0$}    
  \end{subfigure}
  \begin{subfigure}[b]{0.21\linewidth}
    \centering
    \begin{tikzpicture}[dot/.style={draw,shape=circle,
        minimum size=1pt,inner sep=0pt,outer sep=0pt,fill=black}]
      \coordinate[dot,label=210:$0$] (p0) at (0,0);
      \coordinate[dot,label=$1$] (p1) at (60:1);
      \coordinate[dot,label=-30:$2$] (p2) at (1,0);
      \begin{scope}[decoration={markings,
          mark=at position 0.55 with {\arrow{>}}}]
        \draw[postaction={decorate}] (p0) -- (p1) node[midway,left] {$p$};
        \draw[dashed,postaction={decorate}] (p0) -- (p2)
          node[midway,below] {$r$};
        \draw[postaction={decorate}] (p1) -- (p2) node[midway,right] {$q$};
      \end{scope}
    \end{tikzpicture}
    \caption{$\Lambda^2_1$}\label{fig:KanComplex-b}
  \end{subfigure}
  \begin{subfigure}[b]{0.21\linewidth}
    \centering
    \begin{tikzpicture}[dot/.style={draw,shape=circle,
        minimum size=1pt,inner sep=0pt,outer sep=0pt,fill=black}]
      \coordinate[dot,label=210:$0$] (p0) at (0,0);
      \coordinate[dot,label=$1$] (p1) at (60:1);
      \coordinate[dot,label=-30:$2$] (p2) at (1,0);
      \begin{scope}[decoration={markings,
          mark=at position 0.55 with {\arrow{>}}}]
        \draw[dashed,postaction={decorate}] (p0) -- (p1);
        \draw[postaction={decorate}] (p0) -- (p2);
        \draw[postaction={decorate}] (p1) -- (p2);
      \end{scope}
    \end{tikzpicture}
    \caption{$\Lambda^2_2$}
  \end{subfigure}
  \begin{subfigure}[b]{0.21\linewidth}
    \centering
    \begin{tikzpicture}[dot/.style={draw,shape=circle,
        minimum size=1pt,inner sep=0pt,outer sep=0pt,fill=black},
      line join = round, line cap = round,
      x={(0:1cm)},y={(90:1cm)},z={(225:2mm)},scale=0.35]
      \pgfmathsetmacro{\factor}{1/sqrt(2)};
      \coordinate[dot,label=left:$0$] (p0) at (-2,0,-2*\factor);
      \coordinate[dot,label=above:$1$] (p1) at (0,2,2*\factor);
      \coordinate[dot,label=below:$2$] (p2) at (0,-2,2*\factor);
      \coordinate[dot,label=right:$3$] (p3) at (2,0,-2*\factor);
      
      \draw[fill=gray, opacity=.5] (p0.center)--(p1.center)--(p2.center)--cycle;
      \draw[fill=gray, opacity=.5] (p0.center)--(p1.center)--(p3.center)--cycle;
      \draw[fill=gray, opacity=.5] (p0.center)--(p2.center)--(p3.center)--cycle;
      \begin{scope}[decoration={markings,
          mark=at position 0.55 with {\arrow{>}}}]
        \draw[postaction={decorate}] (p0) -- (p1);
        \draw[postaction={decorate}] (p0) -- (p2);
        \draw[postaction={decorate}] (p0) -- (p3);
        \draw[postaction={decorate}] (p1) -- (p2);
        \draw[postaction={decorate}] (p1) -- (p3);
        \draw[postaction={decorate}] (p2) -- (p3);
      \end{scope}
    \end{tikzpicture}
    \caption{$\Lambda^3_0$}
  \end{subfigure}
  \caption{Умова заповнення Кана в розмірностях 2 та 3.}
  \label{fig:KanComplex}
\end{figure}

Наприклад, на Рис.~\ref{fig:KanComplex-b}, якщо нам задано два $1$-симплекси $p$ і $q$ в $X$ зі спільним кінцем, то існує деякий $2$-симплекс, що представляє як композицію $p$ і $q$ (третя грань $r$), так і внутрішність, що відображає той факт, що $r$ є композицією $p$ і $q$.

Зауважимо, що комплекси Кана дають \emph{неалгебраїчне} поняття $\infty$-групоїда: існують композиції та вищі симпліціальні ототожнення, але немає операцій, що виокремлюють конкретну композицію.

Тут ми підходимо до потенційної пастки: ми не можемо сказати, що гомотопічні типи \emph{є} комплексами Кана, оскільки вони мають різні критерії тотожності: у звичайній математичній практиці ми ототожнюємо дві симпліціальні множини, якщо вони ізоморфні (це вже слабше поняття тотожності, ніж те, що надається теорією множин!), тоді як ототожнення між комплексами Кана $X$ і $Y$, \emph{розглянутими як гомотопічні типи}, має бути \emph{гомотопічною еквівалентністю}. 

І це, мабуть, доречний момент, щоб надати теоретико-типовий погляд на відомі гасла Куайна (Quine)~\cite{Quine1969}:\footnote{Друге також обговорювалося з унівалентної перспективи в \cite{Rodin2017,Tsementzis2017}.}
\begin{enumerate}
\item\label{it:Q1} \emph{Бути — означає бути значенням змінної,} і
\item\label{it:Q2} \emph{Немає сутності без тотожності (No entity without identity).}
\end{enumerate}
У~\ref{it:Q1} ми вимагаємо, крім того, щоб усі змінні були типізовані, тому ми скоріше кажемо, що бути $A$ — означає бути значенням змінної типу $A$ і, що більш важливо, \emph{бути — означає бути елементом типу}, а в~\ref{it:Q2} ми не вимагаємо жодного поняття тотожності між сутностями різних типів, але ми вимагаємо, як суттєвої частини задання типу $A$, щоб тип тотожності $x=_Ay$ для $x,y:A$ був осмисленим і правильно виражав засоби ототожнення елементів $A$: \emph{немає типу без типу тотожності.}

Розбіжність між поняттям тотожності між об'єктами моделі (тут — комплексами Кана) та бажаним поняттям тотожності (тут — гомотопічною еквівалентністю) зазвичай вирішується за допомогою \emph{відносних категорій} як інструменту. Відносна категорія складається з категорії, оснащеної широкою підкатегорією \emph{слабких еквівалентностей}. Це часто уточнюється шляхом додавання додаткових властивостей (наприклад, слабкі еквівалентності задовольняють властивості «два з трьох» або «два з шести») або структури, такої як розшарування (fibrations) та/або корозшарування (cofibrations), що добре взаємодіють зі слабкими еквівалентностями. Особливо зручним є поняття \emph{модельної категорії Квіллена} (Quillen model category), яка дійсно містить як розшарування, так і корозшарування на додаток до слабких еквівалентностей, і передбачається, що вона є повною та коповною.

Категорія симпліціальних множин $\sSet$ може бути оснащена структурою модельної категорії Квіллена, в якій фібрантними об'єктами є комплекси Кана (вони також є кофібрантними, оскільки всі об'єкти є кофібрантними), а слабкими еквівалентностями між комплексами Кана є гомотопічні еквівалентності. Категорія топологічних просторів, або точніше, з технічних причин, категорія компактно породжених топологічних просторів, $\cgTop$, аналогічно може бути оснащена модельною структурою Квіллена, в якій кофібрантними об'єктами є хороші простори (технічно — клітинні комплекси; всі об'єкти є фібрантними), а слабкими еквівалентностями між хорошими просторами є гомотопічні еквівалентності.

Квіллен~\cite{Quillen1967} довів, що ці дві модельні категорії породжують еквівалентні \emph{гомотопічні категорії}. З цією метою він ввів поняття (того, що зараз називається) \emph{еквівалентності Квіллена} між модельними категоріями. Для заданого хорошого простору $X$ відповідний сингулярний комплекс Кана $\Pi_\infty(X)$ має як $n$-симплекси неперервні відображення з топологічного $n$-симплекса в $X$, а для заданого комплексу Кана $A$ відповідний простір є \emph{геометричною реалізацією} $\abs A$, що задається склеюванням топологічних симплексів відповідно до відображень граней та виродження в $A$.

Те, що гомотопічні категорії еквівалентні, є першим кроком до отримання того, чого ми насправді хочемо. Насправді ми хотіли б показати, що модельні категорії Квіллена симпліціальних множин і топологічних просторів породжують еквівалентні \emph{гомотопічні типи} (в обох випадках обмежуючись об'єктами, які лежать у фіксованому універсумі Гротендіка). І яке поняття гомотопічного типу ми повинні використовувати тут? Виявляється, це не має значення, але найпростіше побудувати (великі) комплекси Кана з будь-якого з них.

Це досягається шляхом розширення як $\sSet$, так і $\cgTop$ до симпліціально збагачених категорій (остання через побудову сингулярного комплексу Кана на рівні просторів відображень) так, щоб вони стали \emph{симпліціальними модельними категоріями}, а потім взяттям гомотопічно когерентних нервів підкатегорій гомотопічних еквівалентностей між біфібрантними об'єктами, тобто між модельними об'єктами з обох сторін.

Зверніть увагу, що для отримання хорошої теорії гомотопічних типів у класичній постановці нам, здається, також потрібна хороша теорія \inftyone-категорій, тобто категорій, (слабко) збагачених у гомотопічних типах, щоб також отримати хороший контроль над \emph{універсумом} гомотопічних типів, що, звісно, є ще однією назвою, яка передбачає теоретико-типове поняття універсуму, і який складається з гомотопічних типів, які є малими відносно деякого універсуму Гротендіка.

Існує ще одна модельна структура на симпліціальних множинах, біфібрантними об'єктами якої є \emph{квазікатегорії} — ті, що задовольняють послабленню умов заповнення Кана, що робить їх придатними як моделі \inftyone-категорій. Це поняття було введено Бордманом і Фогтом (Boardman and Vogt)~\cite{BoardmanVogt1973}, а отримана теорія \inftyone-категорій була широко вивчена Жуаялем (Joyal)~\cite{Joyal2002} та Лурі (Lurie)~\cite{LurieHTT} (див. також додаток в~\cite{LurieHTT} для детальної інформації про симпліціальні модельні категорії, як обговорювалося вище).

Моя мета у висвітленні цих технічних деталей полягає не лише в тому, щоб пояснити, як визначаються і обробляються гомотопічні типи в теоретико-множинній математиці, але й дати відчуття пов'язаних з цим тонкощів. Знадобилося багато років, щоб дати хороший звіт про те, як працювати з вищою структурою в теоретико-множинній математиці (часто працюючи в $1$-категорному шарі), і досі залишається багато відкритих питань про те, які конструкції та властивості є інваріантними відносно слабких еквівалентностей всередині модельної категорії та відносно еквівалентностей Квіллена між модельними категоріями. Наприклад, лише нещодавно було встановлено, що спряженість Квіллена завжди індукує спряженість між підлеглими квазікатегоріями, а отже, і спряженість представлених \inftyone-категорій~\cite{MazelGee2016}. Інша низка відкритих питань стосується можливості \emph{алгебраїчних} моделей для $\infty$-групоїдів, де композиція, обернені елементи тощо задаються операціями, а не просто припускається їх існування. Цілком можливо, що теорія типів матиме вплив у цій галузі, див., наприклад, пропозицію Брюнері (Brunerie)~\cite[Додаток]{Brunerie2016}.

Тому не повинно викликати подиву те, що досі залишаються відкриті питання про те, як обробляти вищі структури в HoTT/UF, яка є набагато молодшою спробою. До цих питань ми зараз і перейдемо.

\section{Вищі структури в HoTT/UF}
\label{sec:problems}

Коли Воєводський запропонував використовувати теорію типів як основу для математики, він спирався на розуміння того, що вищі структури в математиці \emph{не} завжди природно є об'єктами вищої \emph{категорії}, але вони \emph{завжди} природно є об'єктами вищого \emph{групоїда}.

Серед $\infty$-групоїдів ми знаходимо \emph{зрізані} (truncated) вищі групоїди, ті, чия структура зосереджена у скінченному діапазоні розмірностей. На найнижчому рівні (рівень зрізання $-2$) ми знаходимо стягувані (contractible) типи, ті, які мають лише один елемент з точністю до ототожнень. По-друге, ми маємо висловлювання (propositions). Це типи, всі типи тотожності яких є стягуваними.

Піднімаючись у розмірностях, далі ми знаходимо множини, всі типи тотожності яких є висловлюваннями, і $1$-групоїди, всі типи тотожності яких є множинами, і так далі. Згадаємо з розділу Альтенкірха (Altenkirch), що ці рівні зрізання мають природну формалізацію в HoTT у термінах предикатів
\[
  \hasDimen : \Ntwo \to \Type \to \Prop,
\]
і що ми маємо відповідні типи $n$-зрізаних типів, \trType n.

Не всі типи є зрізаними. $2$-сфера, наприклад, має структуру в усіх розмірностях, тому вона не є $n$-типом для жодного $n$.

$n$-типи пов'язані з універсумом усіх типів, $\Type$, через конструкцію зрізання, яка відображає тип $X$ у найближчий до нього $n$-тип $\trunc X_n$. Існує конструкція $\tr{\blank}_n : X \to
\trunc X_n$, що породжує еквівалентність
\[
  \blank \circ \tr{\blank}_n : (\trunc X_n \to Y) \to (X \to Y)
\]
для будь-якого $n$-типу $Y$.

Коли ми переходимо до обговорення вищої структури, часто саме незрізані типи найважче побудувати. Основна причина полягає в тому, що ми часто можемо будувати зрізані типи зверху вниз, за розмірностями. Щоб побудувати висловлювання, ми можемо просто вказати тип доказу $P$ того, що висловлювання істинне, а потім, якщо необхідно, взяти пропозиційне зрізання $\trunc{P}_{-1}$.

Я хочу підкреслити на цьому етапі, що множини, які ми обговорювали вище (і в розділах Альтенкірха та Аренса-Норта), \emph{не} є множинами теорії множин! Згідно з висловом Куайна, це різні поняття, оскільки вони мають різні поняття тотожності. Давайте тимчасово використаємо індекси для розрізнення, і писатимемо set$_1$ для множини теоретика множин і set$_2$ для множини структураліста/теоретика гомотопій (це також поняття теоретика моделей). Існує навіть третє поняття множини, set$_0$, яке, ймовірно, є більш фундаментальним, ніж set$_1$ або set$_2$, і яке викладається в початковій школі.

З теоретико-типової точки зору, set$_0$ є просто підмножиною фіксованої універсальної set$_2$, $X$. Тобто ми маємо тип
$\Set_0(X) :\equiv \mathcal P(X) :\equiv (X \to \Prop)$,
що представляє булеан (powerset) $X$. Ми маємо елементарне відношення належності, ${\in}_0: X \times \Set_0(X) \to \Prop$,
і дві set$_0$ рівні, якщо вони мають однакові елементи в цьому сенсі.

Це, звісно, не поняття множини теоретика множин, згідно з яким set$_1$ є елементами (а не підмножинами) універсуму міркувань $U$ (який сам є set$_2$), що оснащений відношенням належності ${\in}_1 : U \times U \to \Prop$, яке задовольняє аксіому об'ємності (і бажано багато інших теоретико-множинних аксіом).

Наївною теоретико-множинною надією було б розв'язати рівняння $U = \mathcal P(U)$ (як ототожнення set$_2$, тобто ізоморфізм). Це неможливо через діагональний аргумент Кантора, але його можна наблизити за допомогою кумулятивної ієрархії $V$ — конструкції, яку можна виконати в HoTT через вищий індуктивний тип~\cite[Розділ~10.5]{HoTTBook}. Тут $V$ є великою set$_2$, яка є найменшим розв'язком рівняння $V = \mathcal P_{\mathrm{small}}(V)$, де $\mathcal P_{\mathrm{small}}(V)
:\equiv \Sigma A : \Set. \Sigma f : A \to V. \mathrm{isInjective}(f)$
є типом малих підмножин $V$. Такі set$_1$ можна розглядати як певні цілком впорядковані дерева, і їх вивчення має досить комбінаторний присмак.

Типовим поняттям множини в HoTT/UF є set$_2$, задана $1$-типом $\Set$, і, здається, саме вона найчастіше використовується в математичній практиці поза теорією множин. Наприклад, майже в усіх математичних контекстах кожну множину можна замінити на ізоморфну копію без зміни сенсу чого-небудь. Звичайно, set$_0$ як елементи булеанів $0$-типів/set$_2$ також зустрічаються скрізь у математичній практиці, але для них поняття теоретика множин і структураліста збігаються.

Аналогічно, поняття категорії розпадається на кілька різних понять: я позначатиму як \emph{прекатегорію} поняття, визначене в Розділі~4.4 глави Аренса-Норта (Ahrens-North), і залишу термін \emph{категорія} без прикрас для унівалентної прекатегорії. Дійсно, у більшості категорно-теоретичних контекстів кожну категорію можна замінити на еквівалентну зі збереженням сенсу. Також корисно мати термін \emph{строга категорія} (strict category)~\cite[Розділ~9.6]{HoTTBook} для прекатегорії, тип об'єктів якої є множиною. З точки зору теоретико-множинної математики, $2$-тип категорій виникає зі структури модельної категорії Квіллена на $1$-категорії строгих категорій.

Більшу частину теорії категорій можна формалізувати в HoTT/UF, використовуючи унівалентне визначення категорії. Прекатегорію можна розглядати як категорію з додатковою структурою, а саме оснащену функтором з $\infty$-групоїда. Для строгої категорії цей функтор має як область визначення $0$-вимірний гомотопічний тип. У теоретико-множинних основах автоматично буде так, що кожну категорію можна оснастити такою строгою структурою, але в UF це є додатковим припущенням, ба більше, конструктивно-гомотопічним табу.

\subsection{Аналітичні та синтетичні аспекти HoTT/UF}

Фундаментальна теорія має бути синтетичною в тому сенсі, що вона описує, як конструювати та міркувати з її фундаментальними об'єктами в термінах постульованих правил. Інакше й бути не могло, бо якби вона описувала «фундаментальні» об'єкти в термінах інших, більш фундаментальних об'єктів, і виводила свої правила з їхніх властивостей, вона навряд чи була б фундаментальною.

Гомотопічні теорії типів є синтетичними теоріями $\infty$-групоїдів. Підхід є глибоко \emph{логічним}, де ми розглядаємо \emph{логіку як теорію інваріантів}, як це започаткував Маутнер (Mautner)~\cite{Mautner1946} і пізніше розвинув Тарський (Tarski) у лекції 1966 року~\cite{Tarski1986}. І Маутнер, і Тарський надихалися підходом до геометрії, викладеним в \emph{Ерлангенській програмі} Кляйна~\cite{Klein1872}. Ідея полягає в тому, що логічними поняттями є ті, що інваріантні відносно максимального поняття симетрій універсумів міркувань. Якщо універсумом міркувань є множина, то відповідною групою симетрій є \emph{симетрична група}, що складається з бієкцій множини на себе, але якщо це вищий гомотопічний тип, то це (вища) група \emph{автоморфізмів}, що складається з усіх власних гомотопічних еквівалентностей.

За аналогією з синтетичними теоріями різних понять геометрії (евклідової, афінної, проективної тощо), гомотопічні теорії типів є синтетичними теоріями гомотопічних типів (а теорії множин є синтетичними теоріями set$_1$), пор. також~\cite{Awodey2018}.

Аналітичний аспект полягає в тому, що вся решта математики, всі математичні об'єкти, їхні типи та їхня структура мають бути розвинені в термінах гомотопічних типів. І ключовим критерієм успіху формалізованого поняття є те, що водно задовольняє те, що Аренс і Норт називають \emph{принципом еквівалентності}, і що я пов'язав з висловом Куайна вище: що тип тотожності фіксує передбачуване поняття ототожнень між математичними об'єктами, які ми моделюємо.

Одним із нових аспектів виконання цієї аналітичної роботи в HoTT є те, що під час визначення структурованого об'єкта вже може бути проблемою отримати правильний \emph{тип-носій} (carrier type). У теоретико-множинних основах підійде будь-яка множина-носій правильної потужності, але в HoTT ми більш розбірливі.

Ми дійсно отримуємо деякі переваги від цієї додаткової ретельності. Наприклад, будь-яка конструкція (яка, пам'ятайте, може бути доведенням висловлювання, заселенням множини тощо), яку ми виконуємо на загальній категорії, гарантовано є інваріантною відносно еквівалентності категорій, і ми можемо використовувати правила типів тотожності для транспортування конструкції вздовж будь-якої еквівалентності.

Порівняйте це з ситуацією в теоретико-множинних основах: там ми повинні доводити інваріантність відносно еквівалентності для будь-якої конструкції на типах розмірності більше $0$. Для множин це не є необхідним, оскільки якщо нам задано множину, для якої поняття ототожнення між елементами задається відношенням еквівалентності, ми можемо взяти фактормножину. Таким чином, ситуація в теорії множин є трохи кращою, ніж у теорії типів до HoTT, де конструкція факторизації множин не була загальнодоступною, що призводило до того, що деякі практики називали «пеклом сетоїдів» (setoid hell). Але в теорії множин та сама проблема виникає для будь-якого математичного типу розмірності більше нуля, тому ми можемо припустити, що формалізації на основі теорії множин зіткнуться з «пеклом вищих групоїдів».

\subsection{Деякі конструкції, які здаються неможливими}
\label{sec:negative-problems}

Оскільки доведення висловлювань у HoTT/UF є окремим випадком створення конструкцій, будь-які відкриті наразі проблеми вважаються конструкціями, які ми ще не знаємо, як виконувати.\footnote{Див. \url{https://ncatlab.org/homotopytypetheory/show/open+problems} для актуального списку відкритих проблем.} Але для деяких з них очікується, що складність полягає не лише в тому, що конструкцію складно виконати за допомогою наявних наразі засобів конструювання, а скоріше в нашій гіпотезі, що будуть потрібні абсолютно нові засоби конструювання.

Основним прикладом, на якому я зосереджуся, є \inftyone-категорії та пов'язані з ними поняття (напів)симпліціальних типів. Інтуїтивно, \inftyone-категорія $\calC$ складається з типу об'єктів $\calC_0$, для кожної пари об'єктів $a,b:\calC_0$ — типу морфізмів $\calC_1(a,b)$, операцій для тотожностей і композиції, операцій, що засвідчують закони одиниці та асоціативності, операцій, що засвідчують вищі закони, які вони повинні задовольняти, і так далі \emph{до нескінченності}. Проблема полягає в тому, щоб придумати спосіб задання всіх цих вищих операцій когерентності в одному типі.

Базовим прикладом \inftyone-категорії з точки зору теорії типів є категорія типів $\calS$, яка має як тип об'єктів універсум $\Type$, а як морфізми з $A$ в $B$ — тип функцій з $A$ в $B$, $\calS_1(A,B):\equiv A \to B$. Тут очевидні операції тотожності та композиції задовольняють усі закони та вищі закони \emph{за визначенням}, тому повинно бути особливо легко показати, що $\calS$ є \inftyone-категорією, тобто, якби ми знали, як визначити тип \inftyone-категорій, \inftyoneCat.

Для успіху UF вирішально, щоб ми мали робоче визначення та теорію \inftyone-категорій. Навіть якщо вони не є \emph{фундаментальними} в тому сенсі, що все будується з них, вони все одно є \emph{фундаментальними} в тому сенсі, що вони є ключовим інструментом у розвитку сучасної вищої алгебри, геометрії та топології.

Проблема визначення \inftyone-категорій еквівалентна проблемі визначення типу симпліціальних типів, \sType. \emph{Симпліціальний тип} — це просто функтор $X : \Delta\op\to\calS$, тому якщо ми знаємо, як визначити \inftyoneCat\ і тип функторів між будь-якими двома $\calC,\calD:\inftyoneCat$, тоді ми можемо визначити \sType. З іншого боку, саму \inftyoneCat\ можна визначити як підтип \sType, що складається з \emph{повних типів Сігала} (також званих~\emph{типами Рецка})~\cite{RiehlShulman2017}.

Проблему визначення \inftyoneCat~також можна звести до іншої, здавалося б простішої проблеми, а саме до визначення типу напівсимпліціальних типів, \ssType. \emph{Напівсимпліціальний тип} — це функтор $X:\Delta_+\op\to\calS$, де $\Delta_+$ — підкатегорія $\Delta$ з тими ж об'єктами, але лише ін'єктивними функціями. Проте, $\Delta_+$ є \emph{прямою категорією}, а саме $0$-зрізаною категорією, де відношення «$x$ має нетотожну стрілку в $y$» є цілком впорядкованим відношенням на \emph{множині} об'єктів, тому ми можемо дати більш прямий опис таким чином: Напівсимпліціальний тип $X$ складається з:
\begin{itemize}
\item типу $0$-симплексів $X_0$, і
\item для кожної пари $0$-симплексів $a_0,a_1:X_0$ — типу $1$-симплексів з $a_0$ в $a_1$, $X_1(a_0,a_1)$, і
\item для кожної трійки $0$-симплексів $a_0,a_1,a_2:X_0$ та $1$-симплексів $a_{01}:X_1(a_0,a_1)$, $a_{02}:X_1(a_0,a_2)$ і $a_{12}:X_1(a_1,a_2)$ — типу $2$-симплексів з межею $a_{01},a_{02},a_{12}$,
\item \emph{і так далі \ldots}
\end{itemize}
Знову ми маємо проблему, що, здається, потрібна нескінченно велика кількість даних, але тут виглядає більш правдоподібним, що індуктивний (або коіндуктивний) підхід міг би спрацювати. Тільки ніхто не зрозумів, як це зробити, і під час неформального опитування дослідників HoTT у Варшаві в 2015 році більшість вважала, що це неможливо.

Ми можемо визначити \inftyone-категорії в термінах напівсимпліціальних типів як повні напів-Сігалові типи~\cite{CapriottiKraus2017}. Ця робота була натхненна аналогічною роботою в класичній постановці~\cite{Harpaz2015}. Таким чином, проблеми визначення \inftyone-категорій, симпліціальних типів і напівсимпліціальних типів є еквівалентними, але наразі недосяжними.

Цікаво порівняти випадок \inftyone-категорій з випадком $\infty$-груп: структура \inftyone-категорії на пунктованому зв'язному типі об'єктів — це те саме, що й $\infty$-моноїд (тип якого ми також не знаємо, як побудувати). Але тип $\infty$-груп — це просто тип пунктованих зв'язних типів, причому типом елементів групи є тип тотожності $\Omega A:\equiv(\pt=_A\pt)$, де $\pt:A$ — позначена точка.

Хоча вищевказані проблеми стосуються «великих» типів, існують також проблеми вищої структури, що стосуються «малих» типів. Наприклад, тривимірна сфера як тип $S^3$ повинна мати структуру $\infty$-групи, оскільки вона є гомотопічним типом групи Лі $SU(2)$. Таким чином, ми повинні мати можливість побудувати гомотопічний тип класифікуючого простору $BSU(2)$ з $\Omega BSU(2)=S^3$, але досі ми не змогли цього зробити. (Ми маємо структуру H-простору, що є першим наближенням~\cite{BuchholtzRijke2018}.) Однак у цьому випадку ми очікуємо, що жодні нові засоби конструювання не потрібні.

Очевидним підходом було б побудувати у звичайний спосіб симпліціальну множину, гомотопічним типом якої є $BSU(2)$. Якби ми тоді мали операцію реалізації $\abs\cdot : \sSet \to \Type$, яка перетворює симпліціальну множину на гомотопічний тип, який вона представляє, тоді ми б закінчили. Але таку операцію реалізації саму по собі, здається, неможливо побудувати в елементарній HoTT!

Як останню, важливу, але більш відкриту проблему дозвольте мені згадати проблему розробки метатеорії HoTT всередині HoTT/UF. Вона має дві сторони: одну відносно легку, і одну досить складну. Відносно легка сторона — синтаксична, але навіть тут є труднощі. Ми можемо представити \emph{екстрисивний} (зовнішній) нетипізований синтаксис і відповідні трансформації, знайомі з теорії компіляторів: вибір поверхневого синтаксису, лексичний і синтаксичний аналіз (parsing) цього синтаксису, а потім перевірку типів. Результатом має бути \emph{інтринсивний} (внутрішній) синтаксис, що містить лише добре типізовані елементи. Інтринсивний синтаксис можна змоделювати за допомогою \emph{фактор-індуктивних типів} (quotient inductive types, QIT)~\cite{AltenkirchKaposi2016}, які вже згадувалися в главі Альтенкірха. Головна складність тут полягає в програмній інженерії: як структурувати як інтринсивний, так і екстрисивний синтаксиси, а також трансформації між ними в достатній загальності, щоб охопити всі види теорій типів, які нас цікавлять.

Більш складна сторона — семантична. Ми хочемо визначити інтерпретації інтринсивного синтаксису у внутрішніх моделях, насамперед канонічну модель, де синтаксичні типи $\vdash A$ відображаються в типи $\sem A$, синтаксичні терми $\vdash a:A$ відображаються в терми $\sem a : \sem A$ і так далі. (З теоретико-доказових міркувань ми очікуємо, що зможемо представити інтерпретацію лише локально, наприклад, теорію типів з $n$ універсумами всередині $(n+1)$-го універсуму, або використовуючи сильніші принципи в цільовій теорії типів, ніж у вихідній.) Шульман назвав цю проблему «примусити HoTT з'їсти саму себе»~\cite{Shulman2014}, для якої напрошується термін \emph{аутофагія} (autophagy).

Проблема в тому, що якщо ми використовуємо інтринсивний синтаксис QIT, то все синтаксичне є (гомотопічною) множиною, і правило елімінації дозволить нам елімінувати лише в множини, тоді як для канонічної моделі ми елімінуємо в $\Type$. І якщо ми спробуємо формалізувати інтринсивний синтаксис, використовуючи незрізаний HIT, то, здається, нам знадобиться нескінченно багато шарів когерентності (що нагадує нам про наші проблеми вище з напівсимпліціальними типами та гомотопічною реалізацією симпліціальних множин).

\section{Можливі подальші засоби конструювання}
\label{sec:solutions}

Тепер, коли ми конкретно побачили як спектр конструкцій, які наразі можливі в (елементарній) HoTT, так і деякі видатні проблеми, що здаються недосяжними, давайте підіб'ємо підсумки.

Перший висновок полягає в тому, що нам би дуже хотілося \emph{довести}, що проблема напівсимпліціальних типів не може бути розв'язана в елементарній HoTT. Але припускаючи це, наступним висновком є те, що елементарна HoTT сама по собі є занадто неповною, щоб виконувати свою фундаментальну роль як основа для UF. Необхідно додати подальші засоби конструювання. Але які саме, і як ми вирішуємо, що додавати?

У певному сенсі ситуація є аналогічною до питання нових аксіом для теорії множин, але є дві основні відмінності. По-перше, ми хочемо використовувати теорію типів як \emph{мову програмування}, і це означає, що для будь-якого запропонованого розширення ми повинні вказати, як нові конструкції \emph{обчислюються} при поєднанні з іншими конструкціями теорії типів. Аксіома унівалентності є болючим місцем у цьому відношенні, оскільки надання їй обчислювального змісту було давньою відкритою проблемою. Зараз вона близька до вирішення за допомогою різних \emph{кубічних теорій типів}~\cite{Angiuli-Harper-Wilson2017,Bezem-Coquand-Huber2014,CCHM2018}, але залишається питання, чи моделюють відповідні модельні структури на різних категоріях кубічних множин $\infty$-групоїди (ми знаємо, що \emph{тестові модельні структури} це роблять~\cite{BuchholtzMorehouse2017}). Питання про те, чи можна надати обчислювальний зміст аксіомі зміни розміру висловлювань, залишається повністю відкритим.

По-друге, ми хочемо використовувати HoTT також і в інших моделях, окрім $\infty$-групоїдів. Існує гіпотеза, що елементарну HoTT можна інтерпретувати в будь-якому \inftyone-топосі: точній зліва локалізації категорії функторів $\calC\op\to\calS$ для малої \inftyone-категорії $\calC$. (Надання точного формулювання та актуального статусу цієї гіпотези завело б нас надто далеко; див.~\cite{HoTTWikiTopos}.)

Деякі з найцікавіших цілей задаються \emph{когезивними \inftyone-топосами} (cohesive \inftyone-toposes), об'єкти яких можна розглядати як геометрично структуровані $\infty$-групоїди. (Див. також главу Шрайбера.) Наприклад, у когезивному \inftyone-топосі гладких $\infty$-групоїдів, \SmoothInftyGpd, ми знаходимо всі гладкі многовиди серед $0$-зрізаних об'єктів. І ми, безумовно, хочемо мати можливість міркувати про гладкі \inftyone-категорії, використовуючи HoTT, інтерпретовану в \SmoothInftyGpd. Тому не годиться пропонувати конструкцію типу \inftyoneCat, яка не може бути осмислено виконана в будь-якому \inftyone-топосі.

На противагу цьому, для проблеми інтерпретації теорії типів у внутрішніх моделях, включно з цими когезивними \inftyone-топосами, нам не потрібно вимагати, щоб засоби для цього самі були доступні в довільних моделях. Здається достатнім мати можливість робити це «на верхньому рівні» (at the top level). Але хоч би як ми вирішували цю проблему, це, ймовірно, має бути зроблено обчислювально осмисленими (конструктивними) засобами, щоб ми могли здійснювати доведення за допомогою відображення (reflection) всередині цих моделей. (Звичайно, не всі моделі, що становлять інтерес, можна буде визначити конструктивно.)

Підсумовуючи, ми очікуємо розшарування гомотопічних теорій типів,

\smallskip

\begin{tabular}{@{\textbullet\ }l@{\ :\ }l}
  \HoTTua  & MLTT плюс аксіома унівалентності та зміна розміру висловлювань, \\
  \HoTTel  & \HoTTua\ плюс кодекартові квадрати, \\
  \HoTTelp & \HoTTel\ плюс конструкції, необхідні для \inftyoneCat, \\
  \HoTTuf  & \HoTTelp\ плюс рефлексивні конструкції.
\end{tabular}

\smallskip

\noindent Тут \HoTTua~є основою для спроби формалізації UniMath Воєводського~\cite{Voevodsky2015,UniMath}. Ми знаємо, що \HoTTel~є строго сильнішою (в сенсі наявності меншої кількості моделей; а не в сенсі теоретико-доказової сили), оскільки з \HoTTua~сумісно, що $n$-й універсум є $n$-типом, тоді як якщо універсуми замкнені відносно кодекартових квадратів, то вони не є зрізаними. Оскільки у нас немає доведень неможливості щодо конструкцій \inftyoneCat~або аутофагії (autophagy), все ще можливо, що ми можемо взяти \HoTTelp~та \HoTTuf~як \HoTTel.

Для кожної з них ми можемо розглянути можливість додавання класичних аксіом, таких як закон виключеного третього (LEM) або аксіома вибору (AC). Їх можна розглядати як конструкції, які \emph{ми} загалом не знаємо, як виконувати, але які могла б виконати всезнаюча істота. Я не буду тут ставати на чийсь бік у дебатах про те, чи краще займатися математикою з цими принципами чи без них — скажу лише, що мені здається, що фундаментальна теорія повинна надавати своїм користувачам вибір. (І в будь-якому випадку, ми хочемо теорії без цих конструктивних табу для міркувань про інші топоси, що становлять інтерес.)

Ми також можемо вилучити аксіому зміни розміру, щоб отримати (узагальнені) предикативні системи, і як для предикативних, так і для імпредикативних систем ми можемо додати різні узагальнені індуктивні типи, щоб збільшити теоретико-доказову силу, якщо це необхідно, при цьому зберігаючи системи конструктивними і без зміни їхнього класу моделей \inftyone-топосів. Таким чином, наведене вище розшарування призначене насамперед для розрізнення (передбачуваних) моделей теорій, а не їхньої теоретико-доказової сили, що є ортогональним питанням.

Маючи все це на увазі, давайте обговоримо деякі можливі подальші засоби конструювання, які ми могли б додати як для \HoTTel{}, так і для \HoTTuf.

\subsection{Симпліціальна теорія типів}

Для будь-якого \inftyone-топоса $\calC$ ми можемо розглядати симпліціальні об'єкти в $\calC$, тобто функтори $\Delta\op \to \calC$, і це знову \inftyone-топос. Як згадувалося вище, ми знаходимо в ньому повну підкатегорію \inftyone-категорій відносно $\calC$. Це є основою для пропозиції Ріля та Шульмана (Riehl and Shulman)~\cite{RiehlShulman2017} щодо синтетичної теорії типів для \inftyone-категорій. В їхній теорії типів, назвімо її sHoTT, типи інтерпретуються як \emph{симпліціальні типи}, і вони дають визначення для \emph{типів Сігала} та \emph{типів Рецка}, причому останні представляють \inftyone-категорії. Вони також можуть визначити тип \emph{дискретних} симпліціальних типів, що представляють звичайні типи/$\infty$-групоїди, але цей тип не є типом Рецка, а отже, не є представленням \inftyone-категорії звичайних типів $\calS$. Дійсно, на даний момент неясно, чи це взагалі можна представити в їхній системі.

Попереду ще багато роботи, перш ніж ми зможемо судити про те, наскільки ця теорія типів є корисною для міркувань про \inftyone-категорії. Але з філософської точки зору не може бути задовільним розглядати sHoTT як фундаментальну теорія, яка, наприклад, відіграє роль \HoTTuf. (І, звичайно, вона такою і не задумувалася!) Тому що симпліціальні типи напрошуються на те, щоб їх аналізували саме як такі: симпліціальні об'єкти в категорії типів, а не сприймали як неаналізовані самі по собі. Ми дійсно хочемо мати можливість \emph{визначати} симпліціальні типи всередині теорії, де типи є фундаментальними об'єктами.

\subsection{Дворівневі теорії типів}

Інший підхід до розв'язання проблеми визначення симпліціальних типів полягає в наявності ще одного шару над унівалентною теорією типів, в якому можна міркувати про інфінітарні строгі конструкції, включно з (напів)симпліціальними типами. Однією з пропозицій є Гомотопічна система типів (Homotopy Type System, HTS) Воєводського~\cite{Voevodsky2013}. Ця система розрізняє фібрантні та нефібрантні типи, і вона має два типи тотожності: звичайний гомотопічний тип тотожності, який елімінує лише у фібрантні типи, а також новий нефібрантний тип строгої рівності, що задовольняє правило відображення (reflection rule): якщо $e:a\seq_A a'$ є мешканцем типу строгої рівності, то $a$ і $a'$ є рівними за визначенням. Це правило, звісно, робить перевірку типів нерозв'язною, що вимагає подальшої мови доказів для виведень типів.

Іншою пропозицією є дворівнева теорія типів (2LTT)~\cite{AltenkirchCapriottiKraus2016,AnnenkovCapriottiKraus2017}. Вона подібна до HTS тим, що розрізняє фібрантні та нефібрантні типи (останні називаються \emph{претипами} (pretypes)), але замість правила відображення для типу строгої рівності вона додає правило єдиності доведень рівності (K Штрайхера) та екстенсіональність функцій як аксіому. Таким чином, перевірка типів є розв'язною, але аксіома екстенсіональності функцій ламає обчислення.

Щоб визначити симпліціальні типи в 2LTT, потрібен додатковий принцип поза межами базової конфігурації. Це може бути припущення про те, що фібрантні та нефібрантні натуральні числа збігаються, або більш технічне припущення про те, що фібрантні за Ріді (Reedy) діаграми фібрантних типів, індексовані строгою категорією Ріді, мають фібрантні границі. Перше суттєво обмежує клас доступних моделей, тоді як друге — ні.

Хоч якими б корисними не виявилися дворівневі теорії типів, вони також здаються незадовільними з фундаментальної перспективи. Тому що що таке претип? Претипи можна вмотивувати лише через моделі HoTT, як вони описані в теоретико-множинній математиці, де вони виникають як об'єкти модельної категорії, що представляє \inftyone-топос. Але вони не зберігаються при еквівалентності \inftyone-топосів, тому вони не мають сенсу, незалежного від представлення. Здається, вони є просто зручним інструментом.

\subsection{Обчислювальні теорії типів}

Якщо ми обмежимося однією конструктивною моделлю, то існує принциповий спосіб надання сенсу новим конструкціям. Це робиться через парадигму обчислювальних теорій типів у традиції Nuprl~\cite{Nuprl-book}. Тут ми розглядаємо конкретну модель, щоб дати \emph{пояснення сенсу} (meaning explanations) для суджень теорії типів. Для певного поняття кубічних множин це зробила група Харпера~\cite{Angiuli-Harper-Wilson2017,AngiuliHarper2018}. Перевага такого підходу полягає в тому, що він гарантує, що всі конструкції є обчислювально осмисленими і мають сенс у моделі. Недоліком є те, що він прив'язаний до конкретної моделі (хоча, якщо розсудливо ставитися до того, які примітиви додаються до обчислювальної мови, цей недолік можна мінімізувати), і що він також призводить до теорії типів без розв'язної перевірки типів, тому потрібна окрема теорія доказів або мова свідчень.

\subsection{Аксіоми представлення}

Деяким читачам, можливо, спало на думку, що проблеми, обговорені в розділі~\ref{sec:problems}, слід вирішувати так само, як вони вирішуються в гомотопічній математиці, заснованій на теорії множин, а саме шляхом роботи з представленнями на основі множин.

Ми вже згадували геометричну реалізацію — операцію, яка продукує базовий гомотопічний тип даної симпліціальної множини або топологічного простору. Ми могли б розглянути можливість додавання $\abs\cdot : \sSet \to \Type$ як базової конструкції та аксіоми, що стверджує, що $\abs\cdot$ є сюр'єктивною, що означає, що кожен тип є просто (merely) еквівалентним геометричній реалізації деякої симпліціальної множини. І, можливо, нам слід додати ще одну аксіому, яка б стверджувала, що кожна функція $A \to B$ між типами виникає (просто) як геометрична реалізація функції між репрезентуючими симпліціальними множинами.

Щось подібне дійсно може бути доречним на рівні \HoTTuf, якби йому можна було надати обчислювальний зміст. Але, безумовно, не на рівні \HoTTelp, оскільки аксіоми суттєво обмежили б спектр моделей (вони є конструктивно-гомотопічними табу).

Навіть набагато слабша аксіома \emph{множини покривають} (sets cover, SC), яка стверджує, що кожен тип допускає сюр'єкцію з множини, допускає просту контрмодель — \inftyone-топос~\cite{nLabTypesCover}.

З іншого боку, SC (або щось подібне) необхідна для опису семантики HoTT (з універсумами) в топосах передпучків (presheaf toposes). Дійсно, універсум у топосі передпучків будується з певних \emph{множин}, що покривають $1$-типи передпучків малих множин.

\section{Висновок}
\label{sec:conclusion}

Вищі структури є водночас сенсом існування (raison d'{\^e}tre) і, поки що, ахіллесовою п'ятою HoTT/UF з фундаментальної перспективи. HoTT може з легкістю справлятися з багатьма важливими вищими структурами, такими як $1$-тип множин і $2$-тип категорій, які можуть бути лише недосконало представлені в інших фундаментальних системах. Але досі вона не може визначити (незрізаний) тип \inftyone-категорій, і це є головною перешкодою для фундаментальних прагнень HoTT/UF. Звичайно, HoTT може (і успішно використовується) для міркувань про (структуровані) гомотопічні типи. У такий спосіб різні гомотопічні теорії типів функціонують як \emph{предметно-орієнтовані мови} (domain specific languages, DSLs).

Однак, щоб бути фундаментальною, нам потрібно знайти переконливу конструкцію та теорію \inftyone-категорій, а також семантику HoTT-DSL всередині самої гомотопічної теорії типів. Здається, потрібні нові методи конструювання, але наразі неясно, якими вони мають бути.

Драматична можливість, не згадана в розділі~\ref{sec:solutions}, полягає в тому, що ми все ж повинні вважати \inftyone-категорії фундаментальними і побудувати синтетичну теорію типів, де типами є \inftyone-категорії, а не $\infty$-групоїди. Це була б \emph{спрямована} (directed) теорія типів. Така річ, безсумнівно, була б досить складною через необхідність відстежувати варіантності (variance) (див.~\cite{LicataHarper2011,Nuyts2015} щодо деяких попередніх спроб), і вона представляла б собою повернення до старих способів мислення про категорні основи, хоча й оновлених з урахуванням гомотопічної перспективи. Ми виокремлювали б $\infty$-групоїди як ті \inftyone-категорії, всі морфізми яких є оборотними, замість того, щоб намагатися будувати \inftyone-категорії з $\infty$-групоїдів. У будь-якому випадку, спрямовані теорії типів мають бути корисними також як DSL для міркувань про $(\infty,2)$-топоси.

Особисто я вважаю, що ми знайдемо деяке рішення, яке дозволить нам залишитися на рівні $\infty$-групоїдів для фундаментальної теорії. Можливо, існує різновид дворівневої теорії типів, який дозволяє нам вловити строгу природу \inftyone-категорії типів без постулювання купи безглуздих претипів.

Аналогію, мабуть, можна провести з основами стабільної теорії гомотопій. \inftyone-категорія спектрів є симетричною моноїдальною стабільною \inftyone-категорією, і з фундаментальної точки зору це правильний погляд, оскільки спектри повинні ототожнюватися, коли вони слабо еквівалентні. Однак було виявлено, що ця \inftyone-категорія може бути представлена симетричними моноїдальними модельними категоріями Квіллена (тобто дуже строгими структурами), і це було дуже важливим для полегшення обчислень у стабільній теорії гомотопій~\cite{EKMM1995,MMSS2001}. (Я маю зазначити, що наразі ведеться робота Фінстера, Лікати, Морхауса і Райлі (Finster-Licata-Morehouse-Riley) над розробкою HoTT-DSL для стабільної теорії гомотопій, націленої на когезивний \inftyone-топос параметризованих спектрів: це фіксує строгу моноїдальну структуру спектрів у теорії типів.)

Може бути й так, що для реалізації фундаментального потенціалу HoTT/UF нам знадобиться вловити строгу структуру самої теорії типів, можливо, шляхом відображення більшої кількості структури суджень на рівні типів.

Я впевнений, що хороше рішення зрештою буде знайдено. Ця сфера все ще молода, і буде цікаво побачити, що принесе майбутнє.

\begin{thebibliography}{10}
\providecommand{\url}[1]{{#1}}
\providecommand{\urlprefix}{URL }
\expandafter\ifx\csname urlstyle\endcsname\relax
  \providecommand{\doi}[1]{DOI~\discretionary{}{}{}#1}\else
  \providecommand{\doi}{DOI~\discretionary{}{}{}\begingroup
  \urlstyle{rm}\Url}\fi

\bibitem{Ahrens-Kapulkin-Shulman2015}
Ahrens, B., Kapulkin, K., Shulman, M.: Univalent categories and the {R}ezk
  completion.
\newblock Math. Structures Comput. Sci. \textbf{25}(5), 1010--1039 (2015).
\newblock \doi{10.1017/S0960129514000486}

\bibitem{AltenkirchCapriottiKraus2016}
Altenkirch, T., Capriotti, P., Kraus, N.: Extending homotopy type theory with
  strict equality.
\newblock In: Computer science logic 2016, \emph{LIPIcs. Leibniz Int. Proc.
  Inform.}, vol.~62, pp. Art. No. 21, 17. Schloss Dagstuhl. Leibniz-Zent.
  Inform., Wadern (2016)

\bibitem{AltDanKra2017}
Altenkirch, T., Danielsson, N.A., Kraus, N.: Partiality, revisited.
\newblock In: J.~Esparza, A.S. Murawski (eds.) Foundations of Software Science
  and Computation Structures: 20th International Conference, FOSSACS 2017, Held
  as Part of the European Joint Conferences on Theory and Practice of Software,
  ETAPS 2017, Uppsala, Sweden, April 22-29, 2017, Proceedings, pp. 534--549.
  Springer, Berlin, Heidelberg (2017).
\newblock \doi{10.1007/978-3-662-54458-7_31}

\bibitem{AltenkirchKaposi2016}
Altenkirch, T., Kaposi, A.: Type theory in type theory using quotient inductive
  types.
\newblock In: Proceedings of the 43rd Annual ACM SIGPLAN-SIGACT Symposium on
  Principles of Programming Languages, POPL '16, pp. 18--29. ACM, New York, NY,
  USA (2016).
\newblock \doi{10.1145/2837614.2837638}

\bibitem{AngiuliHarper2018}
Angiuli, C., Harper, R.: Meaning explanations at higher dimension.
\newblock Indagationes Mathematicae \textbf{29}(1), 135--149 (2018).
\newblock \doi{10.1016/j.indag.2017.07.010}.
\newblock L.E.J. Brouwer, fifty years later

\bibitem{Angiuli-Harper-Wilson2017}
Angiuli, C., Harper, R., Wilson, T.: Computational higher-dimensional type
  theory.
\newblock In: POPL '17: Proceedings of the 44th Annual ACM SIGPLAN-SIGACT
  Symposium on Principles of Programming Languages, pp. 680--693. ACM (2017).
\newblock \doi{10.1145/3009837.3009861}

\bibitem{AnnenkovCapriottiKraus2017}
Annenkov, D., Capriotti, P., Kraus, N.: Two-level type theory and applications
  (2017).
\newblock \urlprefix\url{https://arxiv.org/abs/1705.03307}.
\newblock Preprint

\bibitem{Awodey2014}
Awodey, S.: Structuralism, invariance, and univalence.
\newblock Philos. Math. (3) \textbf{22}(1), 1--11 (2014).
\newblock \doi{10.1093/philmat/nkt030}

\bibitem{Awodey2018}
Awodey, S.: Univalence as a principle of logic.
\newblock Indagationes Mathematicae  (2018).
\newblock \doi{10.1016/j.indag.2018.01.011}

\bibitem{Baez2007}
Baez, J.: The homotopy hypothesis (2007).
\newblock \urlprefix\url{http://math.ucr.edu/home/baez/homotopy/}.
\newblock Lecture at Higher Categories and Their Applications, Thematic Program
  on Geometric Applications of Homotopy Theory, Fields Institute, Toronto

\bibitem{Bezem-Coquand-Huber2014}
Bezem, M., Coquand, T., Huber, S.: A model of type theory in cubical sets.
\newblock In: 19th {I}nternational {C}onference on {T}ypes for {P}roofs and
  {P}rograms (TYPES 2013), \emph{LIPIcs. Leibniz Int. Proc. Inform.}, vol.~26,
  pp. 107--128. Schloss Dagstuhl. Leibniz-Zent. Inform., Wadern (2014).
\newblock \doi{10.4230/LIPIcs.TYPES.2013.107}

\bibitem{BoardmanVogt1973}
Boardman, J., Vogt, R.: Homotopy invariant algebraic structures on topological
  spaces, vol. 347.
\newblock Springer, Cham (1973)

\bibitem{Brunerie2016}
Brunerie, G.: On the homotopy groups of spheres in homotopy type theory.
\newblock Ph.D. thesis, Laboratoire J.A. Dieudonn{\'e} (2016).
\newblock \urlprefix\url{https://arxiv.org/abs/1606.05916}

\bibitem{BDR2018}
Buchholtz, U., van Doorn, F., Rijke, E.: Higher groups in homotopy type theory
  (2018).
\newblock \urlprefix\url{https://arxiv.org/abs/1802.04315}.
\newblock Accepted for Proceedings of Logic in Computer Science (LICS 2018)

\bibitem{BuchholtzFavonia2018}
Buchholtz, U., Favonia~(Hou), K.B.: Cellular cohomology in homotopy type theory
  (2018).
\newblock \urlprefix\url{https://arxiv.org/abs/1802.02191}.
\newblock Accepted for Proceedings of Logic in Computer Science (LICS 2018)

\bibitem{BuchholtzMorehouse2017}
Buchholtz, U., Morehouse, E.: Varieties of cubical sets.
\newblock In: P.~Höfner, D.~Pous, G.~Struth (eds.) RAMICS 2017: Relational and
  Algebraic Methods in Computer Science, \emph{Lecture Notes in Computer
  Science}, vol. 10226. Springer, Cham (2017).
\newblock \doi{10.1007/978-3-319-57418-9_5}

\bibitem{BuchholtzRijke2017}
Buchholtz, U., Rijke, E.: The real projective spaces in homotopy type theory.
\newblock In: 32nd Annual ACM/IEEE Symposium on Logic in Computer Science (LICS
  2017), pp. 1--8. IEEE, New York, NY, USA (2017).
\newblock \doi{10.1109/LICS.2017.8005146}

\bibitem{BuchholtzRijke2018}
Buchholtz, U., Rijke, E.: The {Cayley-Dickson} construction in homotopy type
  theory (2018).
\newblock To appear in Higher Structures

\bibitem{CapriottiKraus2017}
Capriotti, P., Kraus, N.: Univalent higher categories via complete semi-{Segal}
  types (2017).
\newblock \urlprefix\url{https://arxiv.org/abs/1707.03693}.
\newblock Preprint

\bibitem{CCHM2018}
Cohen, C., Coquand, T., Huber, S., M{\"o}rtberg, A.: Cubical type theory: a
  constructive interpretation of the univalence axiom.
\newblock In: 21st {I}nternational {C}onference on {T}ypes for {P}roofs and
  {P}rograms (TYPES 2015), LIPIcs. Leibniz Int. Proc. Inform. Schloss Dagstuhl.
  Leibniz-Zent. Inform., Wadern (2018).
\newblock \doi{10.4230/LIPIcs.TYPES.2015.5}

\bibitem{Nuprl-book}
Constable, R.L., Allen, S.F., Bromley, H.M., Cleaveland, W.R., Cremer, J.F.,
  Harper, R.W., Howe, D.J., Knoblock, T.B., Mendler, N.P., Panangaden, P.,
  Sasaki, J.T., Smith, S.F.: Implementing Mathematics with the {Nuprl} Proof
  Development System.
\newblock Prentice-Hall, NJ (1986).
\newblock \urlprefix\url{http://www.nuprl.org/}

\bibitem{vanDoorn2016}
van Doorn, F.: Constructing the propositional truncation using non-recursive
  hits.
\newblock In: Proceedings of the 5th ACM SIGPLAN Conference on Certified
  Programs and Proofs, pp. 122--129 (2016).
\newblock \doi{10.1145/2854065.2854076}

\bibitem{Doorn2018}
van Doorn, F.: On the formalization of higher inductive types and synthetic
  homotopy theory.
\newblock Ph.D. thesis, Carnegie Mellon University (2018)

\bibitem{DRB2017}
van Doorn, F., von Raumer, J., Buchholtz, U.: Homotopy type theory in {Lean}.
\newblock In: M.~Ayala-Rinc{\'o}n, C.~Mu{\~n}oz (eds.) ITP 2017: Interactive
  Theorem Proving, \emph{Lecture Notes in Computer Science}, vol. 10499.
  Springer, Cham (2017).
\newblock \doi{10.1007/978-3-319-66107-0_30}

\bibitem{EKMM1995}
{Elmendorf}, A., {K\v{r}\'{\i}\v{z}}, I., {Mandell}, M.A., {May}, J.: Modern
  foundations for stable homotopy theory.
\newblock In: Handbook of algebraic topology, pp. 213--253. North-Holland,
  Amsterdam (1995)

\bibitem{Grothendieck1983}
Grothendieck, A.: Pursuing stacks (1983).
\newblock \urlprefix\url{http://thescrivener.github.io/PursuingStacks/}.
\newblock Manuscript

\bibitem{Harpaz2015}
Harpaz, Y.: Quasi-unital {$\infty$}-categories.
\newblock Algebraic \& Geometric Topology \textbf{15}(4), 2303–2381 (2015).
\newblock \doi{10.2140/agt.2015.15.2303}

\bibitem{Joyal2002}
Joyal, A.: Quasi-categories and {Kan} complexes.
\newblock Journal of Pure and Applied Algebra \textbf{175}(1-3), 207--222
  (2002).
\newblock \doi{10.1016/S0022-4049(02)00135-4}

\bibitem{KanIII}
{Kan}, D.M.: Abstract homotopy. {III}.
\newblock {Proc. Natl. Acad. Sci. USA} \textbf{42}, 419--421 (1956).
\newblock \doi{10.1073/pnas.42.7.419}

\bibitem{Klein1872}
Klein, F.: {Vergleichende Betrachtungen {\"u}ber neuere geometrische
  Forschungen}.
\newblock Verlag von Andreas Deichert, Erlangen (1872)

\bibitem{Kraus2016}
Kraus, N.: Constructions with non-recursive higher inductive types.
\newblock In: Proceedings of the 31st Annual ACM/IEEE Symposium on Logic in
  Computer Science (LiCS'16), pp. 595--604. ACM, New York, NY, USA (2016).
\newblock \doi{10.1145/2933575.2933586}

\bibitem{KrausSattler2015}
Kraus, N., Sattler, C.: Higher homotopies in a hierarchy of univalent
  universes.
\newblock ACM Trans. Comput. Logic \textbf{16}(2), 18:1--18:12 (2015).
\newblock \doi{10.1145/2729979}

\bibitem{LicataFinster2014}
Licata, D.R., Finster, E.: {Eilenberg-MacLane} spaces in homotopy type theory.
\newblock In: Proceedings of the Joint Meeting of the Twenty-Third EACSL Annual
  Conference on Computer Science Logic (CSL) and the Twenty-Ninth Annual
  ACM/IEEE Symposium on Logic in Computer Science (LICS), CSL-LICS '14, pp.
  66:1--66:9. ACM, New York, NY, USA (2014).
\newblock \doi{10.1145/2603088.2603153}

\bibitem{LicataHarper2011}
Licata, D.R., Harper, R.: 2-dimensional directed type theory.
\newblock Electronic Notes in Theoretical Computer Science \textbf{276},
  263--289 (2011).
\newblock \doi{10.1016/j.entcs.2011.09.026}.
\newblock Twenty-seventh Conference on the Mathematical Foundations of
  Programming Semantics (MFPS XXVII)

\bibitem{LurieHTT}
Lurie, J.: Higher topos theory, \emph{Annals of Mathematics Studies}, vol. 170.
\newblock Princeton University Press, Princeton, NJ (2009).
\newblock \doi{10.1515/9781400830558}

\bibitem{MMSS2001}
{Mandell}, M., {May}, J., {Schwede}, S., {Shipley}, B.: Model categories of
  diagram spectra.
\newblock {Proc. Lond. Math. Soc. (3)} \textbf{82}(2), 441--512 (2001).
\newblock \doi{10.1112/S0024611501012692}

\bibitem{Mautner1946}
Mautner, F.: An extension of {K}lein's {E}rlanger {P}rogram: Logic as
  invariant-theory.
\newblock American Journal of Mathematics \textbf{68}(3), 345--384 (1946).
\newblock \doi{10.2307/2371821}

\bibitem{MazelGee2016}
Mazel-Gee, A.: {Quillen} adjunctions induce adjunctions of quasicategories.
\newblock New York Journal of Mathematics \textbf{22}, 57--93 (2016).
\newblock \urlprefix\url{http://nyjm.albany.edu/j/2016/22-4.html}

\bibitem{Nuyts2015}
Nuyts, A.: Towards a directed homotopy type theory based on 4 kinds of
  variance.
\newblock Ph.D. thesis, KU Leuven (2015).
\newblock
  \urlprefix\url{https://people.cs.kuleuven.be/~dominique.devriese/ThesisAndreasNuyts.pdf}

\bibitem{Quillen1967}
Quillen, D.: Homotopical algebra, vol.~43.
\newblock Springer, Cham (1967).
\newblock \doi{10.1007/BFb0097438}

\bibitem{Quine1969}
Quine, W.V.O.: Ontological relativity and other essays.
\newblock Columbia University Press (1969)

\bibitem{RiehlShulman2017}
Riehl, E., Shulman, M.: A type theory for synthetic {$\infty$}-categories.
\newblock Higher Structures \textbf{1}(1), 147--224 (2017).
\newblock
  \urlprefix\url{https://journals.mq.edu.au/index.php/higher_structures/article/view/36}

\bibitem{Rijke2017}
Rijke, E.: The join construction (2017).
\newblock \urlprefix\url{https://arxiv.org/abs/1701.07538}.
\newblock Preprint

\bibitem{Rijke-Shulman-Spitters2017}
Rijke, E., Shulman, M., Spitters, B.: Modalities in homotopy type theory
  (2017).
\newblock \urlprefix\url{https://arxiv.org/abs/1706.07526}.
\newblock Preprint

\bibitem{Rodin2017}
Rodin, A.: Venus homotopically.
\newblock {IfCoLog} Journal of Logics and their Applications \textbf{4}(4),
  1427--1445 (2017).
\newblock \urlprefix\url{http://collegepublications.co.uk/ifcolog/?00013}.
\newblock {Special Issue Dedicated to the Memory of Grigori Mints. Dov Gabbay
  and Oleg Prosorov (Guest Editors)}

\bibitem{HoTTWikiTopos}
Schreiber, U., Shulman, M.: Model of type theory in an {$(\infty,1)$}-topos.
\newblock
  \urlprefix\url{https://ncatlab.org/homotopytypetheory/show/model+of+type+theory+in+an+%28infinity%2C1%29-topos}.
\newblock Revision 14, Homotopy Type Theory wiki

\bibitem{nLabTypesCover}
Schreiber, U., Shulman, M.: {$n$}-types cover.
\newblock \urlprefix\url{https://ncatlab.org/nlab/show/n-types+cover}.
\newblock Revision 6, nLab

\bibitem{Shulman2014}
Shulman, M.: Homotopy type theory should eat itself (but so far, it’s too big
  to swallow) (2014).
\newblock
  \urlprefix\url{https://homotopytypetheory.org/2014/03/03/hott-should-eat-itself/}.
\newblock Blog post

\bibitem{Shulman2017}
Shulman, M.: Elementary {$(\infty,1)$}-topoi (2017).
\newblock
  \urlprefix\url{https://golem.ph.utexas.edu/category/2017/04/elementary_1topoi.html}.
\newblock Blog post

\bibitem{Tarski1986}
Tarski, A.: What are logical notions?
\newblock History and Philosophy of Logic \textbf{7}(2), 143--154 (1986).
\newblock \doi{10.1080/01445348608837096}.
\newblock Ed. by J. Corcoran

\bibitem{Tsementzis2017}
Tsementzis, D.: Univalent foundations as structuralist foundations.
\newblock Synthese \textbf{194}(9), 3583–3617 (2017).
\newblock \doi{10.1007/s11229-016-1109-x}

\bibitem{HoTTBook}
{Univalent Foundations Program}: Homotopy Type Theory: Univalent Foundations of
  Mathematics.
\newblock \url{http://homotopytypetheory.org/book/}, Institute for Advanced
  Study (2013)

\bibitem{Voevodsky2013}
Voevodsky, V.: A simple type system with two identity types.
\newblock \urlprefix\url{https://ncatlab.org/homotopytypetheory/files/HTS.pdf}.
\newblock Started February 23, 2013. Work in progress.

\bibitem{Voevodsky2014}
Voevodsky, V.: The origins and motivations of univalent foundations.
\newblock The IAS Institute Letter  (2014).
\newblock \urlprefix\url{https://www.ias.edu/ideas/2014/voevodsky-origins}

\bibitem{Voevodsky2015}
Voevodsky, V.: An experimental library of formalized mathematics based on the
  univalent foundations.
\newblock Mathematical Structures in Computer Science \textbf{25}(5),
  1278–1294 (2015).
\newblock \doi{10.1017/S0960129514000577}

\bibitem{UniMath}
Voevodsky, V., Ahrens, B., Grayson, D., et~al.: {UniMath --- a computer-checked
  library of univalent mathematics}.
\newblock {available} at \url{https://github.com/UniMath/UniMath}

\end{thebibliography}

\end{document}
